%%%%%%%%%%%%%%
%
%  第一部分 软件概述
%
%%%%%%%%%%%%%%
\section{软件概述}

\subsection{项目背景}
\par
超声检测（Ultrasonic Testing, UT）是工业无损检测中应用最广泛的手段之一。A-Scan（A型显示）通过“时间-幅度”曲线直观反映回波位置与能量，用于评估材料内部缺陷、壁厚与界面反射情况。
\par
在实际工程现场中，检测人员通常需要结合\textbf{增益}、\textbf{声程}、\textbf{声速}、\textbf{闸门}与\textbf{游标}等功能完成定位与定量，并对可疑反射进行记录与报告输出。为了支持培训演示、算法验证与软件工程化开发，本软件提供 A-Scan 波形的仿真采集、闸门触发、缺陷记录与报告追溯能力。

\subsection{编写目的}
\par
本使用说明书用于指导用户正确安装、启动和使用“\SoftwareName（\Version）”，并明确软件的功能边界与操作流程，主要目的如下：
\begin{enumerate}
	\item 作为软件交付与培训的配套文档，降低上手成本；
	\item 规范常用操作步骤（参数设置、闸门检测、缺陷记录、报告导出等）；
	\item 为软件著作权材料撰写提供可审阅的使用说明依据。
\end{enumerate}

\subsection{适用对象}
\par
本软件主要面向以下用户群体：
\begin{enumerate}
	\item 无损检测人员、质量工程师（用于流程演示与操作练习）；
	\item 检测设备与算法研发人员（用于交互原型验证与算法调试）；
	\item 软件测试与验收人员（用于功能验证与报告追溯演示）。
\end{enumerate}

\subsection{适用范围}
\par
本软件当前为\textbf{仿真框架}，适用范围如下：
\begin{enumerate}
	\item A-Scan 波形的实时显示与交互（深色工业风界面）；
	\item 线性闸门（Start/Width/Level）检测与触发显示；
	\item 双游标（A/B）读数与差值计算；
	\item 缺陷记录、删除/清空管理与 HTML 报告导出；
	\item 声速校准与零点延迟（Probe Delay）校准及文件保存/加载。
\end{enumerate}
\par
说明：本软件未接入真实采集硬件，当前波形为仿真信号。若用于工程验收，需在接入实际硬件与完成系统验证后方可使用。

\subsection{软件特点}
\par
软件具备以下特点：
\begin{enumerate}
	\item 采用 MVC 架构，界面与逻辑分离，便于维护与扩展；
	\item 使用 PyQtGraph 实现高性能绘图，支持实时刷新与交互拖拽；
	\item 闸门与游标遵循现场常见 UT 操作习惯，读数直观；
	\item 缺陷记录可追溯，每条记录可关联波形截图并写入报告；
	\item 支持校准参数持久化，便于不同工件/材料快速切换。
\end{enumerate}

\subsection{运行环境}
\par
推荐运行环境如下：
\begin{enumerate}
	\item 操作系统：Windows 10/11、macOS、Linux（软著材料排版建议使用安装了宋体/新罗马字体的环境）；
	\item Python：3.10 及以上版本；
	\item 第三方库：PyQt5、PyQtGraph、NumPy；
	\item 显示要求：建议分辨率不低于 1366$\times$768。
\end{enumerate}

\subsection{术语与缩略语}
\par
为便于阅读，本文对常用术语作简要说明：
\begin{enumerate}
	\item A-Scan：A型显示，时间-幅度波形；
	\item 增益（dB）：放大信号幅度的参数；
	\item 声程（mm）：设置的最大检测路径或显示范围；
	\item 声速（m/s）：材料中超声传播速度，用于时间与深度换算；
	\item 闸门（Gate）：在时间/深度窗口内进行峰值检测与触发的门控区域；
	\item 门限（Level）：闸门触发阈值，通常以屏高比例表示；
	\item 游标（Cursor）：用于读取波形指定位置的时间/深度与幅度；
	\item Probe Delay：探头/系统零点延迟，用于时间零点校正。
\end{enumerate}
